% Part: computability
% Chapter: computability-theory
% Section: complement-ce

\documentclass[../../../include/open-logic-section]{subfiles}

\begin{document}

\olfileid[fr]{cmp}{thy}{cmp}
\olsection{Les ensembles effectivement énumérables ne sont pas stables par complémentation}

Si $A$ est effectivement énumérable, son complémentaire
$\Complement{A}=\Nat\setminus A$ l'est-il toujours ?
Le théorème et le corollaire suivants répondent
par la négative.

\begin{thm}
\ollabel{thm:ce-comp}
Soit $A$ un ensemble quelconque d'entiers naturels.
Alors $A$ est calculable si et seulement si $A$
et $\Complement{A}$ sont tous deux effectivement
énumérables.
\end{thm}

\begin{proof}
Dans un sens, c'est immédiat : si $A$ est calculable,
$\Complement{A}$ l'est aussi
($\Char{A}=1\tsub\Char{\Complement{A}}$), et
les deux ensembles sont donc effectivement
énumérables.

Réciproquement, supposons que $A$ et
$\Complement{A}$ soient effectivement énumérables.
Soient $A$ le domaine de $\cfind{d}$ et
$\Complement{A}$ le domaine de $\cfind{e}$.
Définissons $h$ par
\[
h(x) = \umin{s}{(T(d,x,s) \lor T(e,x,s))}.
\]
Sur l'entrée~$x$, la fonction $h$ cherche donc un
calcul de $\cfind{d}$ ou de $\cfind{e}$ qui
s'arrête. Si $x\in A$, la première recherche
réussit; si $x\in\Complement{A}$, c'est la
seconde. Ainsi $h$ est totale et calculable.
Pour tout~$x$, on a alors $x\in A$ si et seulement
si $T(d,x,h(x))$, c'est-à-dire si le calcul de
$\cfind{d}$ s'arrête. Comme
$T(d,x,h(x))$ est une relation calculable,
$A$ est calculable.\footnote{Corrections éditoriales : l'indice $d$ désigne le domaine $A$, et $e$ son complémentaire. La source écrit à tort $T(e,x,h(x))$ dans le test d'appartenance à $A$ ; nous employons $T(d,x,h(x))$. Le commentaire utilise aussi les indices $d,e$ fixés dans la preuve, plutôt que les $e,f$ de la source.}
\end{proof}

\begin{explain}
On peut aussi le voir directement : pour décider
si $x\in A$, on cherche en parallèle des calculs
qui s'arrêtent pour $\cfind{d}$ et $\cfind{e}$.
L'un des deux doit être trouvé. Si c'est celui
de $\cfind{d}$, alors $x\in A$; sinon,
$x\in\Complement{A}$.
\end{explain}

\begin{cor}
\ollabel{cor:comp-k}
$\Complement{K_0}$ n'est pas effectivement
énumérable.
\end{cor}

\begin{proof}
Nous savons que $K_0$ est effectivement énumérable
mais non calculable. Si $\Complement{K_0}$ était
également effectivement énumérable, $K_0$ serait
calculable par le \olref{thm:ce-comp}, contrairement
au \olref[ncp]{thm:K-0}.
\end{proof}

\end{document}
